\documentclass[aps,prb,reprint,superscriptaddress,nofootinbib]{revtex4-2}

\usepackage[utf8]{inputenc}
\usepackage[T1]{fontenc}
\usepackage{amsmath,amssymb,amsfonts,bm}
\usepackage{mathtools}
\usepackage{braket}
\usepackage{xcolor}
\usepackage{hyperref}

\hypersetup{colorlinks=true,linkcolor=blue,citecolor=blue,urlcolor=blue}

\newcommand{\bG}{\mathbf G}
\newcommand{\bq}{\mathbf q}
\newcommand{\btau}{\boldsymbol\tau}
\newcommand{\ee}{\mathrm e}
\newcommand{\ii}{\mathrm i}
\newcommand{\cell}{\Omega}
\newcommand{\dd}{\mathrm d}

\begin{document}

\title{Note on the preservation of separability under
species-resolved truncated SVD of the local PAW auxiliary}

\author{Internal note}
\affiliation{Companion to \emph{paw\_isdf\_thc\_prb.tex}, Sec.~V.A--V.D.}
\date{\today}

\begin{abstract}
This note clarifies a point raised about the derivation in Sec.~V of the
PAW-ISDF-THC manuscript: how the separability of the local auxiliary
functions
\(\eta_{a\lambda}^{\bq}(\bG)
=\ee^{-\ii(\bq+\bG)\cdot\btau_a}\,\eta_{s(a)}(\lambda;\bq+\bG)\)
is preserved when the species-resolved factor is rotated to a truncated
SVD basis.  The conclusion is that separability is preserved by
construction, because the rotation acts on an index that is absent from
the atom phase, and the only approximation introduced by truncation is a
rank reduction in the channel index.  We also distinguish this exact
algebraic statement from the empirical claim that the \(\bq=0\) basis
also covers \(\eta_s\) at \(\bq\neq 0\), which is a separate, evidence-based
assertion.
\end{abstract}

\maketitle

\section{What ``separability'' means in the original definition}

The starting point is the species-resolved factorization of the local
auxiliary, Eq.~(eta-species-factor) of the manuscript,
\begin{equation}
\eta_{a\lambda}^{\bq}(\bG)
=\ee^{-\ii(\bq+\bG)\cdot\btau_a}\,
\eta_{s(a)}(\lambda;\bq+\bG).
\label{eq:def}
\end{equation}
This is \emph{not} a tensor-product decomposition over the indices
\((\lambda,\bG)\): both factors carry a \(\bG\) dependence.  The
``separability'' that the manuscript claims is structural and concerns
the role of the atom index \(a\):
\begin{itemize}
\item the atom dependence is carried entirely by the unitary phase
\(\ee^{-\ii(\bq+\bG)\cdot\btau_a}\), which has \emph{no} channel index
\(\lambda\);
\item the channel/species dependence is carried entirely by
\(\eta_s(\lambda;\bq+\bG)\), which depends on the atom only through
\(a\!\mapsto\!s(a)\).
\end{itemize}
This particular separation -- atom phase factored out from a species
function -- is what enables a per-species (rather than per-atom) low-rank
compression in the first place: a basis built once for species \(s\)
applies, modulo the phase, to all atoms of that species.

\section{The rotation acts on \texorpdfstring{$\lambda$}{lambda} only}

The truncated singular vectors are stored in a fixed matrix
\(B_s^{(\varepsilon)}\in\mathbb C^{n_\lambda\times k_s}\), of which the
relevant fact is its index structure: it carries an old channel index
\(\lambda\) and a reduced channel index \(\widetilde\lambda\), and
\emph{nothing else}.  Concretely, \(B_s^{(\varepsilon)}\) does not
depend on
\begin{enumerate}
\item the atom \(a\) (only on the species \(s\));
\item the transferred momentum \(\bq\);
\item the reciprocal vector \(\bG\).
\end{enumerate}
This is the entire content of the separability claim under truncation.
Because \(B_s^{(\varepsilon)}\) lives in a different index space from
the atom phase and from the K-argument, contraction with
\(B_s^{(\varepsilon)\,\dagger}\) commutes with both:
\begin{align}
&\sum_\lambda\bigl[B_s^{(\varepsilon)}\bigr]^{*}_{\lambda\widetilde\lambda}
\Bigl[\,\ee^{-\ii(\bq+\bG)\cdot\btau_a}\,
\eta_s(\lambda;\bq+\bG)\Bigr] \notag\\
&\quad=\;\ee^{-\ii(\bq+\bG)\cdot\btau_a}\,
\underbrace{\sum_\lambda
\bigl[B_s^{(\varepsilon)}\bigr]^{*}_{\lambda\widetilde\lambda}\,
\eta_s(\lambda;\bq+\bG)}_{\Phi^{(s)}_{\widetilde\lambda}(\bq+\bG)}.
\label{eq:rotation-commutes}
\end{align}
The right-hand side has the same shape as Eq.~\eqref{eq:def}: an atom
phase multiplying an atom-independent K-space function indexed by the
reduced channel \(\widetilde\lambda\).  This is precisely the definition
of the reduced auxiliary in Eq.~(eta-reduced-spectral),
\begin{equation}
\eta'_{a\widetilde\lambda}(\bG;\bq)
=\ee^{-\ii(\bq+\bG)\cdot\btau_a}\,
\Phi^{(s)}_{\widetilde\lambda}(\bq+\bG).
\label{eq:reduced}
\end{equation}
No assumption about the SVD itself was used in
Eq.~\eqref{eq:rotation-commutes}.  For \emph{any} fixed rectangular
matrix \(B\in\mathbb C^{n_\lambda\times k}\), the same identity holds
and the same separability follows.  The role of the SVD is only to
choose a particular \(B\) that captures most of the spectrum of
\(\eta_s\) at small \(k\).

\section{Where the approximation actually enters}

The expansion (eta-low-rank-expansion) of the manuscript states
\begin{equation}
\eta_{a\lambda}^{\bq}(\bG)
\;\approx\;
\sum_{\widetilde\lambda}
\bigl[B_s^{(\varepsilon)}\bigr]_{\lambda\widetilde\lambda}\,
\eta'_{a\widetilde\lambda}(\bG;\bq).
\label{eq:expansion}
\end{equation}
Substituting Eq.~\eqref{eq:reduced} into Eq.~\eqref{eq:expansion} and
collecting,
\begin{align}
\sum_{\widetilde\lambda}
\bigl[B_s^{(\varepsilon)}\bigr]_{\lambda\widetilde\lambda}\,
\eta'_{a\widetilde\lambda}(\bG;\bq)
&=\ee^{-\ii(\bq+\bG)\cdot\btau_a}\,
\bigl[B_s^{(\varepsilon)}{B_s^{(\varepsilon)}}^{\dagger}
\eta_s(\cdot;\bq+\bG)\bigr]_\lambda.
\label{eq:projector-form}
\end{align}
Comparing this with Eq.~\eqref{eq:def} shows that the
``\(\approx\)'' in Eq.~\eqref{eq:expansion} reduces to a single
operation: replacing the full \(\eta_s(\lambda;\bq+\bG)\) by its
projection onto the column space of \(B_s^{(\varepsilon)}\).  Two
features are notable:
\begin{enumerate}
\item the atom phase is identically preserved -- it factors out of
Eq.~\eqref{eq:projector-form} with no approximation, exactly as in the
original definition;
\item the projector \(B_s^{(\varepsilon)}{B_s^{(\varepsilon)}}^{\dagger}\)
acts on the channel index \(\lambda\) only, not on \(a\), \(\bq\), or
\(\bG\).
\end{enumerate}
In particular, the truncation does not introduce any new dependence on
the atom or on \(\bG\) into the species factor.  Whatever separation
of indices was present in Eq.~\eqref{eq:def} is preserved verbatim by
Eqs.~\eqref{eq:reduced} and \eqref{eq:projector-form}.

\section{Algebraic vs.\ empirical content}

It is helpful to keep two layers of the construction distinct, since
they are sometimes conflated.

\paragraph{Exact (algebraic).}  For \emph{any} \(n_\lambda\times k\)
matrix \(B\), the form
\(\ee^{-\ii(\bq+\bG)\cdot\btau_a}\,\Phi^{(s)}_{\widetilde\lambda}(\bq+\bG)\)
with
\(\Phi^{(s)}_{\widetilde\lambda}(\mathbf K)
=\sum_\lambda B^{*}_{\lambda\widetilde\lambda}\,\eta_s(\lambda;\mathbf K)\)
has the same atom-phase / species-resolved separability as the original
\(\eta_{a\lambda}^{\bq}(\bG)\).  This is the algebraic content of
Sec.~V.A and is independent of how \(B\) was chosen.

\paragraph{Empirical (numerical).}  The specific basis
\(B_s^{(\varepsilon)}\) constructed from a Coulomb-weighted SVD of
\(\eta_s\) at \(\bq=0\) is observed to cover \(\eta_s(\lambda;\bq+\bG)\)
at all \(\bq\neq 0\) with essentially the same projection residual.
Equivalently, the column space of \(B_s^{(\varepsilon)}\) is closed
under the \(\bq\)-shift in argument.  This is the carry-to-q claim of
Sec.~V.D.  The structural reason is that \(\bq\) enters only through
the magnitude \(|\bq+\bG|\) in the radial Bessel transforms and through
the angular factors \(Y_{LM}(\widehat{\bq+\bG})\), both of which vary
smoothly within the same finite-dimensional angular--radial subspace as
\(\bq\) sweeps the IBZ.  The numerical evidence in Sec.~V.D (full
\(2\times 2\times 2\) IBZ for Si PAW/USPP and LiH PAW; residual at
\(\bq\neq 0\) within numerical precision of the residual at \(\bq=0\))
substantiates this empirically.

The two questions are therefore:
\begin{itemize}
\item \emph{Does truncation preserve separability?}  Yes,
unconditionally and by construction
(Eqs.~\eqref{eq:rotation-commutes}--\eqref{eq:projector-form}).
\item \emph{Does the \(\bq=0\) truncation also approximate \(\eta_s\) at
\(\bq\neq 0\) well?}  This is an empirical statement, not an algebraic
one, and is the content of the carry-to-q test in Sec.~V.D.
\end{itemize}
A concern about ``separability after truncation'' is therefore really
either trivial (the algebraic case) or about the rank captured by the
\(\bq=0\) SVD at finite \(\bq\) (the empirical case); these are
qualitatively different and should be evaluated separately.

\section{Separability of the augmentation charge}
\label{sec:Qhat}

The same separability question can be asked of the augmentation pair
density \(\widehat Q^{\,a}_{\alpha\beta}\) that drives the local-ISDF
construction in Eq.~(local-isdf-qhat) of the manuscript.  We show that
the atom-phase / species-resolved separability of \(\widehat Q\)
survives the channel-index truncation by the same mechanism as for
\(\eta\), and discuss the additional structural feature
(rank-1 in the channel-pair indices) that makes the THC ansatz hold.

\subsection{The original separable form}

The augmentation pair density of atom \(a\), species \(s=s(a)\), in
Bloch form on the dense \(\bG\)-grid, is
\begin{equation}
\widehat Q^{\,a}_{\alpha\beta}(\bG;\bq)
=\ee^{-\ii(\bq+\bG)\cdot\btau_a}\,
q^{(s)}_{\alpha\beta}(\bq+\bG),
\label{eq:Qhat-species}
\end{equation}
i.e.\ an atom-dependent unitary phase times the species-resolved
augmentation charge \(q^{(s)}_{\alpha\beta}(\mathbf K)\) (the QE
\(Q^{IJ}(\mathbf K)\) before structure factor).  This is the
\(\widehat Q\) analog of Eq.~(eta-species-factor) and asserts the
same atom-phase / species separation: the atom enters only through the
phase and does not act on the channel-pair indices \((\alpha,\beta)\).

The full-rank local-ISDF in Eq.~(local-isdf-qhat) writes the species
augmentation charge as a sum of rank-1 outer products in the
channel-pair indices,
\begin{equation}
q^{(s)}_{\alpha\beta}(\mathbf K)
=\sum_{\lambda=1}^{n_\lambda}
U^{(s)\,*}_{\lambda\alpha}\,
U^{(s)}_{\lambda\beta}\,
\eta_s(\lambda;\mathbf K),
\label{eq:Qhat-full-rank}
\end{equation}
in the symmetric-pair basis defined by the
\texttt{build\_local\_isdf\_full\_rank} construction
(\(n_\lambda=n_h^2\); each diagonal pair \(I=J\) contributes one row,
each off-diagonal pair \(I<J\) contributes two rows
\((\lambda_+,\lambda_-)\) with \(U_{\lambda_\pm,I}=+1/\sqrt2\),
\(U_{\lambda_+,J}=+1/\sqrt2\), \(U_{\lambda_-,J}=-1/\sqrt2\), and
sign \(\pm1\) absorbed into \(\eta_s(\lambda;\cdot)\)).  At full rank
this representation is exact.

\subsection{Atom-phase separability is preserved unconditionally}

Combining Eqs.~\eqref{eq:Qhat-species}--\eqref{eq:Qhat-full-rank},
\begin{equation}
\widehat Q^{\,a}_{\alpha\beta}(\bG;\bq)
=\ee^{-\ii(\bq+\bG)\cdot\btau_a}
\sum_\lambda U^{(s)\,*}_{\lambda\alpha}\,U^{(s)}_{\lambda\beta}\,
\eta_s(\lambda;\bq+\bG).
\label{eq:Qhat-original}
\end{equation}
The atom enters only through the prefactor; the sum on the right-hand
side has no \(a\)-dependence beyond \(a\!\mapsto\!s(a)\).  Now apply
the truncation argument of Sec.~II to the species function under the
sum.  Replacing \(\eta_s(\lambda;\mathbf K)\) by its rank-\(k_s\)
projection gives, with \(\Pi_s:=B_s^{(\varepsilon)}{B_s^{(\varepsilon)}}^\dagger\),
\begin{align}
\widehat Q^{\,a}_{\alpha\beta}(\bG;\bq)
&\approx
\ee^{-\ii(\bq+\bG)\cdot\btau_a}
\sum_\lambda U^{(s)\,*}_{\lambda\alpha}\,U^{(s)}_{\lambda\beta}\,
\bigl[\Pi_s\,\eta_s(\cdot;\bq+\bG)\bigr]_\lambda. \notag\\
&\quad
\label{eq:Qhat-projected}
\end{align}
The atom phase factors out identically: the truncation operator
\(\Pi_s\) acts on the channel index \(\lambda\) alone and has no
\(a\)-, \(\bq\)-, or \(\bG\)-dependence.  Thus the species-resolved
factorization \eqref{eq:Qhat-species} survives -- the truncated
species charge is still atom-independent, the K-argument still
\(\bq+\bG\), and the atom dependence still lives entirely in the phase.
This is the augmentation-charge analog of the result in
Sec.~II--III for \(\eta\).

\subsection{The rank-1 channel-pair structure}

The \(\widehat Q\) factorization in Eq.~\eqref{eq:Qhat-full-rank} has
an additional structural property that the THC ansatz of
Eq.~(final-paw-isdf-thc) requires: the coefficient at each channel row
\(\lambda\) is a \emph{rank-1 outer product} in the channel-pair
indices, \(U^{(s)\,*}_{\lambda\alpha}U^{(s)}_{\lambda\beta}\).  This is
what permits the linear-in-projector feature
\(Y_{a\lambda,i}=\sum_\alpha U_{a,\lambda\alpha}P_{i,a\alpha}\) of
Eq.~(Y-feature), and hence the
\(X_{\Lambda i}^{*}X_{\Lambda j}\) form of the smooth-grid pair-density
factorization in Eq.~(comp-rho-local-isdf).

Two compression strategies are visible in the present codebase and
manuscript.  They differ in how the rank-1 structure interacts with
the channel-index truncation, and it is worth being explicit about
this distinction.

\subsubsection{Strategy A -- pair-subset selection (current implementation)}

The implementation in \texttt{local\_isdf\_compress.hpp} performs a
pivoted Cholesky on the pair-space Gram matrix in the \(L^2\) or
Coulomb metric and \emph{selects a subset of the original symmetric
pairs}.  Each kept pair contributes its original symmetric-pair rows
(one diagonal row, or a \(\pm\) split) verbatim from the full-rank
construction.  The compressed \(U\) is \(U\) restricted to those rows;
the compressed \(\eta_s\) is \(\eta_s\) restricted to the same rows.
The augmentation-charge representation in this case is
\begin{equation}
q^{(s)}_{\alpha\beta}(\mathbf K)
\approx
\sum_{\lambda\in\mathcal P_s}
U^{(s)\,*}_{\lambda\alpha}\,U^{(s)}_{\lambda\beta}\,
\eta_s(\lambda;\mathbf K),
\label{eq:Qhat-subset}
\end{equation}
where \(\mathcal P_s\) is the kept pair set.  The rank-1 outer-product
structure in \((\alpha,\beta)\) is preserved trivially -- the truncation
discards entire \(\lambda\) rows, and what remains is a direct subset
of the original sum.  The \(K_a\) closed form
(\texttt{local\_isdf.hpp::compute\_K\_a}) inherits this structure: the
round-trip identity
\(\sum_\lambda \mathrm{sign}(\lambda)U_{\lambda I}U_{\lambda J} f(i(\lambda),j(\lambda))=f(I,J)\)
holds restricted to \((I,J)\in\mathcal P_s\), and \(\Delta C\)
contributions from dropped pairs are reconstructed as zero.
Atom-phase separability and rank-1 channel-pair structure are both
exact, with truncation error governed by the pivoted-Cholesky residual
in the chosen metric.

\subsubsection{Strategy B -- channel-index SVD rotation
(manuscript Sec.~V.B): numerical refutation}

The manuscript Sec.~V.B describes a different compression: rotate the
channel index by the truncated SVD basis,
\(U'_{a,\widetilde\lambda\alpha}
=\sum_\lambda[B_s^{(\varepsilon)}]^{*}_{\lambda\widetilde\lambda}
U_{a,\lambda\alpha}\) (Eq.~(U-reduced)), and
\(\Phi^{(s)}_{\widetilde\lambda}(\mathbf K)
=\sum_\lambda[B_s^{(\varepsilon)}]^{*}_{\lambda\widetilde\lambda}
\eta_s(\lambda;\mathbf K)\) (Eq.~(phi-species)).  The implicit claim is
that the augmentation charge can then be reconstructed in the rotated
basis as
\begin{equation}
q^{(s),\,B}_{\alpha\beta}(\mathbf K)
:=\sum_{\widetilde\lambda}
U'^{(s)\,*}_{\widetilde\lambda\alpha}\,
U'^{(s)}_{\widetilde\lambda\beta}\,
\Phi^{(s)}_{\widetilde\lambda}(\mathbf K),
\label{eq:Qhat-Bs-claim}
\end{equation}
parallel to Eq.~\eqref{eq:Qhat-full-rank}.  This claim is
\emph{numerically false}, and not by a small amount.  Substituting
the rotation definitions on the right-hand side gives
\begin{equation}
q^{(s),\,B}_{\alpha\beta}(\mathbf K)
=\sum_{\lambda_1\lambda_2\lambda_3}
T_{\lambda_1\lambda_2\lambda_3}\,
U^{(s)}_{\lambda_1\alpha}\,U^{(s)}_{\lambda_2\beta}\,
\eta_s(\lambda_3;\mathbf K),
\label{eq:Qhat-Bs-trilinear}
\end{equation}
with the trilinear kernel
\begin{equation}
T_{\lambda_1\lambda_2\lambda_3}
=\sum_{\widetilde\lambda}
\bigl[B_s^{(\varepsilon)}\bigr]_{\lambda_1\widetilde\lambda}\,
\bigl[B_s^{(\varepsilon)}\bigr]^{*}_{\lambda_2\widetilde\lambda}\,
\bigl[B_s^{(\varepsilon)}\bigr]^{*}_{\lambda_3\widetilde\lambda}.
\label{eq:trilinear-T}
\end{equation}
For Eq.~\eqref{eq:Qhat-Bs-claim} to reduce to the original full-rank
identity Eq.~\eqref{eq:Qhat-full-rank}, one would need
\(T_{\lambda_1\lambda_2\lambda_3}
=\delta_{\lambda_1\lambda_2}\delta_{\lambda_1\lambda_3}\),
which holds for \(B_s=I\) but \emph{not} for any other unitary -- in
particular not for the SVD basis of \(\eta_s\).

The diagnostic test \texttt{paw\_qhat\_strategy\_b}
(\texttt{src/hamiltonian/tests/test\_hamilt.cpp}) builds
\(B_s^{(\varepsilon)}\) from the full Coulomb-weighted SVD of
\(\eta_s\) at \(\bq=0\), reconstructs
\(q^{(s),B}_{\alpha\beta}(\mathbf K)\) per
Eq.~\eqref{eq:Qhat-Bs-claim}, and compares it to the
\(\mathrm{qgm}_{\alpha\beta}(\mathbf K)\) stored in the pseudopotential.
Reporting the relative Coulomb-norm residual
\(\bigl\|q^{(s),B}-\mathrm{qgm}\bigr\|_C/\bigl\|\mathrm{qgm}\bigr\|_C\):
\begin{table}[h]
\centering
\begin{tabular}{lccccc}
\hline
species & \(n_\lambda\) & full-rank id. & \(k=n_\lambda\) (no trunc.) & \(k(10^{-3})\) & \(k(10^{-9})\) \\
\hline
LiH Li PAW & 64  & \(4{\times}10^{-19}\) & \(\mathbf{1.97}\) & \(\mathbf{1.97}\) & \(\mathbf{1.97}\) \\
LiH H  PAW & 4   & \(1{\times}10^{-20}\) & \(\mathbf{2.61}\) & \(\mathbf{2.61}\) & \(\mathbf{2.61}\) \\
Si    PAW  & 324 & \(5{\times}10^{-20}\) & \(\mathbf{1.93}\) & \(\mathbf{1.93}\) & \(\mathbf{1.93}\) \\
\hline
\end{tabular}
\end{table}

The middle column verifies that
\(\sum_\lambda U_{\lambda I}U_{\lambda J}\eta_\lambda=\mathrm{qgm}_{IJ}\)
holds at machine precision in the symmetric-pair basis -- the original
full-rank ISDF identity is exact.  The remaining three columns show
that the Strategy~B reconstruction is off by an O(1) factor at every
truncation rank, including the full rank \(k=n_\lambda\) where the
SVD basis itself introduces \emph{no} approximation.  The error is the
trilinear obstruction Eq.~\eqref{eq:trilinear-T}, not an SVD truncation
residual: the SVD basis cannot make
\(\sum_\lambda B B^{*}B^{*}\) into \(\delta\delta\), and the
manuscript's substitution
\(\widehat Q\to\sum_{\widetilde\lambda}U'^{*}U'\Phi\) is therefore
\emph{wrong as a representation of the augmentation charge}.

The same conclusion applies to \(\Delta C_{a,\alpha\beta,\gamma\delta}\)
in Eq.~(dC-reduced) of the manuscript: the substitution of
\(U^{*}U\to U'^{*}U'\) and \(K\to K'=B^\dagger K B\) into
Eq.~(local-k-factorization) does not reproduce the original
\(\Delta C\) under generic \(B_s\), by the same trilinear obstruction
acting on each pair of \((\alpha,\beta)\) and \((\gamma,\delta)\)
indices.

\subsubsection{Summary}

Atom-phase separability of \(\widehat Q\) is preserved unconditionally
by either compression strategy, by the argument in Secs.~II--III: the
truncation acts on an index that is absent from both the atom phase
and the K-argument, so the factorization
\(\widehat Q^{\,a}=\ee^{-\ii(\bq+\bG)\cdot\btau_a}\,q^{(s)}\) survives
unchanged and the truncated \(q^{(s)}\) remains atom-independent.

The rank-1 outer-product structure in \((\alpha,\beta)\), however,
\emph{cannot} be carried through a generic SVD rotation:
\(\widetilde C^{(s)}_{\widetilde\lambda,\alpha\beta}
=\sum_\lambda U^{*}_{\lambda\alpha}U_{\lambda\beta}B_{\lambda\widetilde\lambda}\)
is a weighted sum of rank-1 outer products and is generically of
higher rank in \((\alpha,\beta)\).  Replacing \(\widetilde C^{(s)}\)
by \(U'^{(s)\,*}U'^{(s)}\) (Strategy~B as written in the manuscript)
introduces an O(1) error, demonstrated numerically above.

The implementation in \texttt{local\_isdf\_compress.hpp} avoids this
problem by selecting a \emph{subset} of the original symmetric-pair
rows (Strategy~A): the kept rows retain the rank-1 form verbatim, the
roundtrip identity holds restricted to kept pairs, and the truncation
error is governed by the pivoted-Cholesky residual in the chosen
metric -- a single, well-controlled approximation.  This is the
strategy that should be referenced in Sec.~V.B; the SVD-rotation
language as currently written in the manuscript is incorrect for
\(\widehat Q\) and \(\Delta C\), even though it is correct for the
species function \(\eta_s\) on its own (Secs.~II--III of this note).

\section{Numerical search for tighter compressions}
\label{sec:tighter}

Strategy A preserves the rank-1 form by construction but is restricted
to a \emph{subset} of the original \(n_\lambda=n_h^2\) channel rows.
A natural question is whether a different rank-1 outer-product
decomposition can do better -- i.e.\ whether
\begin{equation}
q^{(s)}_{\alpha\beta}(\mathbf K)\;\approx\;
\sum_{k=1}^{K} u_{k\alpha}\,u_{k\beta}\,v_k(\mathbf K),
\qquad K \ll n_h^2,
\label{eq:cp-symmetric}
\end{equation}
admits a small \(K\) for given Coulomb-norm tolerance.  Equation
\eqref{eq:cp-symmetric} is the symmetric tensor-rank (CP) decomposition
of \(q^{(s)}_{\alpha\beta}(\mathbf K)\) viewed as a complex symmetric
3-tensor in \((\alpha,\beta,\mathbf K)\), and finding the minimal \(K\)
is in general NP-hard.  Two heuristic constructions were attempted
(\texttt{paw\_qhat\_greedy\_rank1} and
\texttt{paw\_qhat\_outer\_svd\_inner\_takagi} in
\texttt{src/hamiltonian/tests/test\_hamilt.cpp}); both fall short.

\subsection{Greedy rank-1 ALS}

The greedy approach extracts a sequence of rank-1 components by
power-iteration on Takagi (best symmetric rank-1 fit), subtracts each
from the residual, and repeats:
\begin{enumerate}
\item Initialize \(u\in\mathbb C^{n_h}\) (unit norm).
\item \(v_g=\sum_{\alpha\beta}\bar u_\alpha\bar u_\beta R_{\alpha\beta g}\).
\item \(M_{\alpha\gamma}=\sum_g w_C(g)\,v_g^{*}\,R_{\alpha\gamma g}\).
\item \(u\to M\bar u/\|M\bar u\|\).  Iterate until converged.
\item \(R\gets R-u\,u^{T}\otimes v\); next \(k\).
\end{enumerate}
For Si PAW (\(n_h=18\), \(n_\lambda=324\)) the greedy residual is
\(8.6\times 10^{-2}\) at \(K=100\), and never reaches \(10^{-3}\)
within the \(K\le 100\) budget.  For LiH Li PAW (\(n_h=8\),
\(n_\lambda=64\)), \(K=64\) reaches only \(5\times 10^{-4}\) -- well
short of machine precision, even though the original symmetric-pair
construction is exact at the same \(K\).  Greedy CP is famously
trapped by local minima; the leading rank-1 components are correct,
but later ones must fit a deflated residual that is no longer
well-aligned with any single rank-1 outer product.

\subsection{Outer SVD + inner Takagi}

A constructive 2-level decomposition truncates first in the
\((\alpha\beta,g)\) Coulomb-weighted unfolding, then Takagi-decomposes
each leading symmetric \(n_h\times n_h\) component:
\begin{equation}
q^{(s)}_{\alpha\beta}(g)\;\approx\;
\sum_{k=1}^{K_{\text{out}}}\sigma_k\,W_k(\alpha,\beta)\,V_k(g),
\quad
W_k=\sum_{m=1}^{r_k}\mu_m^{(k)}w_m^{(k)}{w_m^{(k)}}^{T},
\label{eq:outer-inner}
\end{equation}
giving a total of \(\sum_k r_k\) rank-1 outer-product terms.  In the
test on Si PAW, the leading \(W_k\) matrices are essentially full
inner rank (typical \(r_k(10^{-6})\!\sim\!12\text{--}18\) out of
\(n_h=18\)), so the 2-level construction \emph{inflates} the
auxiliary dimension well above \(n_\lambda\):

\begin{table}[h]
\centering
\begin{tabular}{lccc}
\hline
\(\varepsilon_{\text{out}}\) & inner \(\varepsilon=10^{-2}\) & \(10^{-3}\) & \(10^{-6}\) \\
\hline
LiH Li (\(n_\lambda=64\), Coul) \\
\quad \(10^{-3}\) & 57 & 73  & 86 \\
\quad \(10^{-6}\) & 134 & 166 & 186 \\
\quad \(10^{-9}\) & 214 & 246 & 266 \\
Si PAW (\(n_\lambda=324\), Coul) \\
\quad \(10^{-3}\) & 619 & 684 & 842 \\
\quad \(10^{-6}\) & 1033 & 1186 & 1419 \\
\quad \(10^{-9}\) & 2415 & 2614 & 2981 \\
\hline
\end{tabular}
\end{table}

Even the loosest setting (\(\varepsilon_{\text{out}}=10^{-3}\),
\(\varepsilon_{\text{in}}=10^{-2}\)) gives \(619>n_\lambda=324\) for
Si.  The leading symmetric directions in the
\((\alpha\beta,g)\) unfolding are simply not aligned with single rank-1
outer products: each \(W_k\) is a generic complex symmetric matrix
that requires \(\mathcal O(n_h)\) Takagi modes to represent.

\subsection{What this means}

These experiments suggest that the symmetric tensor rank of
\(q^{(s)}_{\alpha\beta}(\mathbf K)\) is essentially tight at
\(n_\lambda=n_h^2\) under the rank-1-per-term constraint required by
the THC ansatz.  The \(\eta_s\) Coulomb-rank diagnostic of
Sec.~V.A--V.D in the manuscript (Si: \(k_\eta\sim 100\) at
\(\varepsilon=10^{-6}\)) measures the rank of the species function
\(\eta_s(\lambda;\mathbf K)\) as a 2-tensor in \((\lambda,\mathbf K)\).
That rank is small.  But coupling \(\eta_s\) to the symmetric-pair
structure factor \(U_{\lambda\alpha}U_{\lambda\beta}\) requires
\(n_h^2\) rank-1 outer products in \((\alpha,\beta)\), one per
\(\lambda\), and the empirical evidence is that no other rank-1 basis
of similar size does any better.

The practical conclusion: Strategy A (pivoted Cholesky pair-subset on
the symmetric-pair Gram in the chosen metric, as already implemented
in \texttt{local\_isdf\_compress.hpp}) is the right route.  It
respects the rank-1-per-term structure exactly, achieves the
information-theoretic lower bound for that ansatz class, and is
controlled by a single tolerance.  The Phase 4.5 plan should be
updated to describe Strategy A explicitly and drop the
\(B_s\)-rotation language entirely.

If genuinely smaller compression is required, three escape routes are
worth flagging, none free:
\begin{enumerate}
\item Replace the symmetric form \(u_{k\alpha}u_{k\beta}v_k(g)\) by a
non-symmetric form \(p_{k\alpha}q_{k\beta}v_k(g)\) (different left and
right factors).  This breaks the THC \(X=Y^{*}Y\) ansatz and forces a
generalized THC kernel.
\item Allow simultaneous (non-greedy) ALS on
Eq.~\eqref{eq:cp-symmetric} with multiple random restarts.  This is
expensive and offers no rank guarantees -- but the symmetric tensor
rank may be smaller than the upper bounds the simple heuristics
above provide.
\item Exploit physical structure further: angular-momentum coupling
selection rules (Wigner 3-j) restrict which \((\alpha,\beta)\) pairs
appear at given \((L,M)\) channels of \(\eta_s\), suggesting a
block-diagonal structure that Strategy A does not exploit.
\end{enumerate}

\section{The local Coulomb tensor \texorpdfstring{$\Delta C$}{Delta C}: same story}
\label{sec:deltaC}

The same compression question can be asked of the one-center correction
\(\Delta C_{a,\alpha\beta,\gamma\delta}\), which the manuscript factors
in the same channel basis as
\begin{equation}
\Delta C_{a,\alpha\beta,\gamma\delta}
=\sum_{\lambda\xi}
U_{\lambda\alpha}\,U_{\lambda\beta}\,
K_{a,\lambda\xi}\,
U_{\xi\gamma}\,U_{\xi\delta},
\label{eq:dC-K}
\end{equation}
with \(K_a\) the \(n_\lambda\times n_\lambda\) real symmetric Coulomb
kernel.  Two perspectives on compressibility, again testing the THC
form against an unconstrained matrix-rank bound.

\subsection{Matrix rank of \texorpdfstring{$K_a$}{K_a}}

The eigendecomposition \(K_a=\sum_k\sigma_k v_k v_k^{T}\) gives the
optimal rank-\(r\) approximation of the unfolding
\(\Delta C_{(\alpha\beta),(\gamma\delta)}\).  Because \(\Delta C\) is
the difference between AE and PS Coulomb tensors,
\(K_a\) is generically \emph{indefinite} (signed eigenvalues).  Test
\texttt{paw\_K\_a\_compressibility}:

\begin{table}[h]
\centering
\begin{tabular}{lccccc}
\hline
species & \(n_\lambda\) & \(r(10^{-3})\) & \(r(10^{-6})\) & \(r(10^{-9})\) & signs (+/-/0) \\
\hline
LiH Li & 64  & \textbf{5}   & 15 & 23  & 16 / 17 / 31 \\
Si PAW & 324 & \textbf{60}  & 97 & 113 & 87 / 76 / 161 \\
\hline
\end{tabular}
\end{table}

For Si PAW, half the eigenvalues are zero (\(K_a\) has effective rank
\(\le 163\)), and at relative tolerance \(10^{-3}\) the matrix rank is
\(60\) -- a \(5.4\times\) compression of the unfolding.

\subsection{Constrained ALS: rank-1 form
\texorpdfstring{$\sum u_\alpha u_\beta v_\gamma v_\delta$}{Sum u u v v}}

To compress \(\Delta C\) while keeping rank-1 outer products in each
pair (the THC-preserving constraint, analogous to
Eq.~\eqref{eq:cp-symmetric} for \(q^{(s)}\)) we ran greedy rank-1 ALS
on
\begin{equation}
\Delta C_{\alpha\beta\gamma\delta}
\;\approx\;
\sum_{k=1}^{K}\,u_{k\alpha}u_{k\beta}\,v_{k\gamma}v_{k\delta}.
\label{eq:dC-uuvv}
\end{equation}
Test \texttt{paw\_deltaC\_greedy\_rank1}.  Convergence is much slower
than the unconstrained matrix-rank bound suggests:

\begin{table}[h]
\centering
\begin{tabular}{lccc}
\hline
species & \(n_\lambda\) & ALS rank @ \(10^{-3}\) & residual @ \(K=200\) \\
\hline
LiH Li & 64  & 28      & --   \\
Si PAW & 324 & \emph{not reached} & \(7\times 10^{-2}\) \\
\hline
\end{tabular}
\end{table}

For Si PAW, even \(K=200\) terms in the constrained ALS form reach
only \(\sim7\times 10^{-2}\) relative Frobenius residual -- compared
to the unconstrained matrix rank \(60\) at \(10^{-3}\) and rank \(97\)
at \(10^{-6}\) from the eigendecomposition.

\subsection{Why the constrained form costs so much more}

The eigendecomposition of \(K_a\) writes
\(\Delta C\!=\!\sum_k\sigma_k W_k(\alpha,\beta)W_k(\gamma,\delta)\)
with each \(W_k(\alpha,\beta)=\sum_\lambda v_{k\lambda}U_{\lambda\alpha}U_{\lambda\beta}\)
a generic symmetric \(n_h\times n_h\) matrix -- not necessarily
rank-1.  Forcing every term into the form
\(u_\alpha u_\beta\) instead requires further Takagi-decomposing each
\(W_k\) into up to \(n_h\) rank-1 outer products, multiplying the
total auxiliary count by a factor of order \(n_h\).  The outer-SVD +
inner-Takagi obstruction of Sec.~\ref{sec:tighter} reappears here in
exactly the same form.

\subsection{Conclusion: same trilinear obstruction, different tensor}

The local Coulomb tensor exhibits the same pattern as the augmentation
charge:
\begin{itemize}
\item As a generic matrix in \((\alpha\beta,\gamma\delta)\) the
unfolding is genuinely low rank (\(60\) of \(324\) for Si at
\(10^{-3}\)).
\item Under the rank-1-per-term constraint required by the THC
\(X^{*}\!X\) ansatz, that low rank cannot be exploited -- greedy ALS
plateaus and the constrained tensor rank is much higher.
\item Using the matrix-rank compression directly (rotating to the
\(K_a\) eigenbasis and identifying \(\widetilde\lambda\to k\)) breaks
the THC ansatz by exactly the same trilinear obstruction analyzed in
Sec.~\ref{sec:Qhat} for \(\widehat Q\): the auxiliary-rotated
expression \(\sum_k\widetilde U_k\widetilde U_k\,(\sigma_k)\,\widetilde U_k\widetilde U_k\)
involves
\(\sum_\lambda v_{k\lambda}v_{k\lambda'}v_{k\lambda''}\), which does
not reduce to a Kronecker-delta product.
\end{itemize}
For Phase 4.5 of the implementation, this leaves Strategy A (pivoted
Cholesky pair-subset on the symmetric-pair Gram) as the only
THC-respecting compression of either \(\widehat Q\) or \(\Delta C\).
A genuinely lower-dimensional THC representation of the local block
would require giving up the \(X^{*}X\) ansatz on the local rows,
e.g.\ a generalized factorization
\(\sum_k X^Y_{ki}X^Z_{kj}\) with two distinct features per local row.

\section{Redefining the augmentation by moment-matching only}
\label{sec:moment-matching}

A different way out of the trilinear wall is to relax the augmentation
function itself.  The PAW augmentation
\(\widehat Q^{\,a}_{\alpha\beta}(\mathbf r)\) is constrained
\emph{only} by its multipole moments
\begin{equation}
q^{LM}_{\alpha\beta}
\;=\;\int\!\dd^3 r\; r^L Y_{LM}(\hat r)
\bigl[\rho^{a,\mathrm{AE}}_{\alpha\beta}(\mathbf r)
-\rho^{a,\mathrm{PS}}_{\alpha\beta}(\mathbf r)\bigr],
\label{eq:moments}
\end{equation}
which is a \emph{finite} set of linear conditions; the radial shape of
\(\widehat Q\) is otherwise free.  QE's canonical construction
\begin{equation}
\widehat Q^{\,a}_{\alpha\beta}(\mathbf r)
=\sum_{LM}a^{LM}_{\alpha\beta}\,
q^{l_\alpha l_\beta}_L(|\mathbf r|)\,Y_{LM}(\hat r)
\label{eq:QE-aug}
\end{equation}
is one such choice, with \(a^{LM}_{\alpha\beta}\) the angular Wigner
coupling and \(q^{l_\alpha l_\beta}_L(r)\) the radial Q function.

If one postulates instead a strictly separable form
\begin{equation}
\widehat Q^{\,a}_{\alpha\beta}(\mathbf r)
\;\approx\;\sum_{k=1}^{K}\,
u_{k\alpha}\,u_{k\beta}\,f_k(\mathbf r),
\label{eq:aug-CP}
\end{equation}
with \(f_k\) any localized shape, the moment-matching condition
becomes a constraint on the moments \(\mu_k^{LM}\) of \(f_k\):
\begin{equation}
q^{LM}_{\alpha\beta}
\;=\;\sum_k u_{k\alpha}u_{k\beta}\,\mu_k^{LM}.
\label{eq:moment-match-CP}
\end{equation}
Since \(f_k(\mathbf r)\) is free up to its moments, the
\(\mu_k^{LM}\) are free parameters; the entire problem reduces to a
symmetric CP decomposition of the \emph{moment tensor}
\(q^{LM}_{\alpha\beta}\), which lives on a tiny finite-dimensional
auxiliary index \((LM)\) -- not the dense G-grid where the trilinear
obstruction was so painful.

\subsection{Numerical test on the moment tensor}

Test \texttt{paw\_moment\_separable} builds
\(q^{LM}_{\alpha\beta}=a^{LM}_{\alpha\beta}\cdot\mathrm{augmom}(L,\beta_\alpha,\beta_\beta)\)
from QE's \texttt{ap} table and per-species \texttt{augmom} radial
moments, then reports two ranks: (i) the \emph{matrix unfolding} rank
in \((\alpha\beta\text{-pair},LM)\), which is the rank achievable in
the more flexible Tucker form, and (ii) the symmetric CP rank from
greedy rank-1 ALS of Eq.~\eqref{eq:aug-CP}.

\begin{table}[h]
\centering
\begin{tabular}{lcccccc}
\hline
species & \(n_h\) & \(n_\lambda\) & \(n_{LM}\) & matrix \(@10^{-6}\) & matrix \(@10^{-9}\) & ALS \(@10^{-3}\) \\
\hline
LiH Li & 8  & 64  & 9  & 9   & 25  & 34 \\
LiH H  & 2  & 4   & 9  & 1   & 3   & 1  \\
Si PAW & 18 & 324 & 25 & \textbf{25}  & 130 & not reached \\
\hline
\end{tabular}
\end{table}

The matrix unfolding rank is the rank achievable in the
\emph{Tucker} form
\begin{equation}
q^{LM}_{\alpha\beta}
\;=\;\sum_{k=1}^{K_{\text{Tuck}}}\,
W_k(\alpha,\beta)\,\nu_k^{LM},
\label{eq:aug-Tucker}
\end{equation}
where each \(W_k\) is a generic symmetric matrix in \((\alpha,\beta)\)
-- not rank-1.  For Si this drops the auxiliary count from
\(n_\lambda=324\) to \(\mathbf{25}\) at \(10^{-6}\) tolerance, a 13×
reduction.  This is the lower bound on what any rank-1-respecting
construction can achieve, and it is genuinely small.

The greedy CP rank, however, does not reach this lower bound -- for
Si, \(K=200\) terms in Eq.~\eqref{eq:aug-CP} still leave
\(\sim4\times 10^{-2}\) relative residual.  Putting each \(W_k\) into
rank-1 outer-product form via Takagi would multiply
\(K_{\text{Tuck}}=25\) by an inner Takagi rank that is empirically
\(\mathcal O(n_h)\); the total is no smaller than \(n_\lambda\), and
the trilinear wall reappears.

\subsection{What this implies for THC}

The conceptual answer to ``can we redefine \(\widehat Q\) as separable
and still match moments'' is unambiguously \textbf{yes}: the moment
condition is finite and linear, and Eq.~\eqref{eq:moment-match-CP}
admits a CP solution at sufficiently large \(K\).  The numerical
answer about how small \(K\) gets divides into two regimes:
\begin{enumerate}
\item \textbf{Tucker-like form} (Eq.~\eqref{eq:aug-Tucker} with
generic \(W_k\)) reaches \(K_{\text{Tuck}}=25\) for Si at
\(10^{-6}\), \(13\times\) below \(n_\lambda=324\).  This form is
\emph{not} compatible with the standard \(X=Y^{*}Y\) THC ansatz:
each \(W_k\) is not rank-1 in \((\alpha,\beta)\).
\item \textbf{Strict CP form} (Eq.~\eqref{eq:aug-CP}, rank-1
outer product per term) is what THC requires.  Greedy ALS for Si
cannot reach \(10^{-3}\) within \(K\le 200\) terms; the symmetric
CP rank of the moment tensor is empirically much larger than the
matrix-unfolding rank, by the same mechanism that defeats every
other rank-1-preserving compression we tried.
\end{enumerate}

In short: the moment-matching reformulation does \emph{not} circumvent
the trilinear obstruction.  It does, however, sharpen the conclusion:
the bottleneck is the rank-1-in-\((\alpha,\beta)\) constraint required
by THC, not the radial-shape freedom that the moment-only condition
provides.  The intrinsic rank of the moment tensor (25 for Si) is
already genuinely small; what fails is the further insistence that
each contribution carry a rank-1 outer product.  Any future low-rank
local-block compression for THC would need to relax the
\(X=Y^{*}Y\) ansatz on the local rows, not the choice of augmentation
function.

\subsection{Asymmetric form: \texorpdfstring{$\sum u_\alpha v_\beta \mu_{LM}$}{Sum u v mu} with \texorpdfstring{$u\neq v$}{u != v}}

A natural generalization is to drop the \((\alpha\!\leftrightarrow\!\beta)\)
symmetry of each rank-1 component:
\begin{equation}
\widehat Q^{\,a}_{\alpha\beta}(\mathbf r)
\;\approx\;\sum_{k=1}^K u_{k\alpha}\,v_{k\beta}\,f_k(\mathbf r),
\quad u_k\neq v_k,
\label{eq:aug-CP-asym}
\end{equation}
which gives the moment-matching condition
\(q^{LM}_{\alpha\beta}=\sum_k u_{k\alpha} v_{k\beta} \mu_k^{LM}\) -- a
plain (asymmetric) CP decomposition of the 3-tensor.  This form has
strictly more freedom than the symmetric one and matches the form
underlying the ``two distinct features per local row''
(\(X^Y_{ki}X^Z_{kj}\)) generalization of THC.

Test \texttt{paw\_moment\_asymmetric} runs greedy rank-1 ALS for
Eq.~\eqref{eq:aug-CP-asym}.  Comparison of residuals at matched \(K\):

\begin{table}[h]
\centering
\begin{tabular}{lccc}
\hline
species & form & residual @ \(K=96\) & residual @ \(K=200\) \\
\hline
Si PAW & symmetric \(uu\)        & \(4.22\times10^{-2}\) & --                   \\
Si PAW & asymmetric \(uv\)       & \(4.55\times10^{-2}\) & \(7.9\times 10^{-3}\) \\
\hline
\end{tabular}
\end{table}

The two forms behave nearly identically for the first \(\sim\!100\)
terms; the asymmetric form pulls ahead modestly at higher \(K\),
reaching \(\sim 8\times 10^{-3}\) at \(K=200\) where the symmetric
form is still at \(\sim 4\times 10^{-2}\).  At neither setting does
the asymmetric CP reach the matrix-unfolding lower bound of \(25\) (Si
at \(10^{-6}\)) or even \(10^{-3}\) within \(K\le 200\).

The structural reason is the same as the outer-SVD-+-inner-Takagi
analysis for \(\widehat Q^{(s)}\) and \(\Delta C\): the leading
matrix-unfolding singular components \(W_k(\alpha,\beta)\) have
SVD-rank \(\sim n_h\) in \((\alpha,\beta)\), so converting them to
rank-1 outer products \(u_{k\alpha}v_{k\beta}\) -- whether symmetric
or asymmetric -- still costs \(\mathcal O(n_h)\) terms per leading
\(W_k\).  Going asymmetric saves a factor of \(\le 2\) (since real
symmetric eigendecomposition pairs eigenvalues vs.\ a single
asymmetric SVD), but does not unlock the matrix-rank compression.

The bottom line is the same as before: the rank-1-on-each-pair-side
constraint is the hard wall, not the symmetry constraint.  Going
asymmetric helps a little; it does not help enough to make local-block
compression below \(n_\lambda\) feasible while keeping a strict
THC-like factorization on each local row.

\section{Suggested clarification in the manuscript}

To prevent the present source of confusion in a referee report, one
sentence inserted immediately after Eq.~(eta-reduced-spectral) is
sufficient.  A possible wording:
\begin{quote}
\itshape
We emphasize that \(B_s^{(\varepsilon)}\) is a fixed
\(n_\lambda\times k_s\) matrix with no dependence on the atom \(a\), the
transferred momentum \(\bq\), or the reciprocal vector \(\bG\).  The
rotation therefore acts on the channel index \(\lambda\) alone and
commutes with both the atom phase and the K-argument.  The separable
structure of Eq.~(eta-species-factor) is preserved by construction, and
the only approximation in Eq.~(eta-low-rank-expansion) is the rank-\(k_s\)
truncation in the channel index.
\end{quote}

\end{document}
